\documentclass[10pt,a4paper]{amsart}
\usepackage{amsmath,amssymb,amsthm,graphicx}
% Retained minimal preamble for the Stacks Project chapter more-algebra.tex (More on Algebra).
% The chapter file pinned-original-source.tex carries only \input{preamble} on its first line, so no
% preamble text lives inside the chapter itself.  The Stacks Project uses one single shared preamble for
% every chapter; the snapshot stages it identically as sources/stacks-chow/preamble.tex and
% sources/stacks-etale-cohomology/preamble.tex.  The block below is copied mechanically, line for line,
% from sources/stacks-chow/preamble.tex lines 14-18 (the Xy-pic setup) and 154-217 (the native theorem
% declarations together with the mathematical macros).  Nothing was added, removed or redefined; the
% wrapper class and the amsmath/amssymb/amsthm/graphicx base packages are supplied by the assembler.
% Chapter: more-algebra.tex -- More on Algebra.
% For commutative diagrams we use Xy-pic
\usepackage[all]{xy}

% We use 2cell for 2-commutative diagrams.
\xyoption{2cell}

% Theorem environments.
%
\theoremstyle{plain}
\newtheorem{theorem}[subsection]{Theorem}
\newtheorem{proposition}[subsection]{Proposition}
\newtheorem{lemma}[subsection]{Lemma}

\theoremstyle{definition}
\newtheorem{definition}[subsection]{Definition}
\newtheorem{example}[subsection]{Example}
\newtheorem{exercise}[subsection]{Exercise}
\newtheorem{situation}[subsection]{Situation}

\theoremstyle{remark}
\newtheorem{remark}[subsection]{Remark}
\newtheorem{remarks}[subsection]{Remarks}

\numberwithin{equation}{subsection}

% Macros
%
\def\lim{\mathop{\mathrm{lim}}\nolimits}
\def\colim{\mathop{\mathrm{colim}}\nolimits}
\def\Spec{\mathop{\mathrm{Spec}}}
\def\Hom{\mathop{\mathrm{Hom}}\nolimits}
\def\Ext{\mathop{\mathrm{Ext}}\nolimits}
\def\SheafHom{\mathop{\mathcal{H}\!\mathit{om}}\nolimits}
\def\SheafExt{\mathop{\mathcal{E}\!\mathit{xt}}\nolimits}
\def\Sch{\mathit{Sch}}
\def\Mor{\mathop{\mathrm{Mor}}\nolimits}
\def\Ob{\mathop{\mathrm{Ob}}\nolimits}
\def\Sh{\mathop{\mathit{Sh}}\nolimits}
\def\NL{\mathop{N\!L}\nolimits}
\def\CH{\mathop{\mathrm{CH}}\nolimits}
\def\proetale{{pro\text{-}\acute{e}tale}}
\def\etale{{\acute{e}tale}}
\def\QCoh{\mathit{QCoh}}
\def\Ker{\mathop{\mathrm{Ker}}}
\def\Im{\mathop{\mathrm{Im}}}
\def\Coker{\mathop{\mathrm{Coker}}}
\def\Coim{\mathop{\mathrm{Coim}}}

% Boxtimes
%
\DeclareMathSymbol{\boxtimes}{\mathbin}{AMSa}{"02}

%
% Macros for moduli stacks/spaces
%
\def\QCohstack{\mathcal{QC}\!\mathit{oh}}
\def\Cohstack{\mathcal{C}\!\mathit{oh}}
\def\Spacesstack{\mathcal{S}\!\mathit{paces}}
\def\Quotfunctor{\mathrm{Quot}}
\def\Hilbfunctor{\mathrm{Hilb}}
\def\Curvesstack{\mathcal{C}\!\mathit{urves}}
\def\Polarizedstack{\mathcal{P}\!\mathit{olarized}}
\def\Complexesstack{\mathcal{C}\!\mathit{omplexes}}
% \Pic is the operator that assigns to X its picard group, usage \Pic(X)
% \Picardstack_{X/B} denotes the Picard stack of X over B
% \Picardfunctor_{X/B} denotes the Picard functor of X over B
\def\Pic{\mathop{\mathrm{Pic}}\nolimits}
\def\Picardstack{\mathcal{P}\!\mathit{ic}}
\def\Picardfunctor{\mathrm{Pic}}
\def\Deformationcategory{\mathcal{D}\!\mathit{ef}}

% --- appended 2026-10-12 (Chief): the project own definitions of the reference/slogan/history
% display environments, copied mechanically from the shared Stacks preamble lines 8-12
% (\usepackage{verbatim} plus the three \newenvironment lines, which the chapter preamble keeps
% verbatim from the project preamble). Nothing added, removed or redefined: these environments
% carry no mathematical content and the project itself typesets them as comments. Without these
% three lines a counted statement whose own text uses them fails with Environment slogan undefined.
\usepackage{verbatim}
\newenvironment{reference}{\comment}{\endcomment}
\newenvironment{slogan}{\comment}{\endcomment}
\newenvironment{history}{\comment}{\endcomment}
\usepackage{cleveref}
\crefname{originallemma}{original lemma}{original lemmas}
\Crefname{originallemma}{Original lemma}{Original lemmas}
\makeatletter
\newlabel{lemma-localize-relative-pseudo-coherent}{{0679}{source}{}{sourcepointer-1}{}}
\newlabel{lemma-localize-relative-pseudo-coherent@cref}{{[originallemma][0][]0679}{[1][0][]source}{}{}{}}
\makeatother
\crefname{originallemma}{original lemma}{original lemmas}
\Crefname{originallemma}{Original lemma}{Original lemmas}
\makeatletter
\newlabel{lemma-finite-push-pseudo-coherent}{{064Z}{source}{}{sourcepointer-2}{}}
\newlabel{lemma-finite-push-pseudo-coherent@cref}{{[originallemma][0][]064Z}{[1][0][]source}{}{}{}}
\makeatother
\crefname{originallemma}{original lemma}{original lemmas}
\Crefname{originallemma}{Original lemma}{Original lemmas}
\makeatletter
\newlabel{lemma-glue-pseudo-coherent}{{066D}{source}{}{sourcepointer-3}{}}
\newlabel{lemma-glue-pseudo-coherent@cref}{{[originallemma][0][]066D}{[1][0][]source}{}{}{}}
\makeatother
\title{Relative pseudo-coherence is local for a cover by principal opens of the algebra}
\author{The Stacks Project Authors}
\date{stacks-project:more-algebra.tex:local-snapshot-20261012; GNU FDL 1.2 or later}
\begin{document}
\maketitle
\setcounter{subsection}{1}
\noindent\textit{Editorial note.} All delivered mathematical text is the author's own, copied from the pinned Stacks Project chapter file more-algebra.tex; no statement or proof step was rewritten, reordered or removed. The counted statement is the labeled lemma environment at source lines 21668-21681 (canonical tag 067E) and the counted proof is the proof environment at source lines 21683-21712, which immediately follows it and is its own proof.\par\medskip
\noindent\textit{Editorial note.} The printed proof begins a sentence with a lower case choose at source line 21686 and writes such that the elements map as described; the wording is copied exactly as printed and is not corrected.\par\medskip
\noindent\textit{Editorial note.} The three references are all to earlier results of this chapter, namely the localization statement, the transfer of relative pseudo-coherence and the glueing statement; the kinds recorded are the type words appearing in the labels themselves. All three are navigation stubs with their roles recorded in the reasons.\par\medskip
\noindent\textit{Editorial note.} The preamble span runs to line 90, the current end of the retained preamble, which includes the appended environment definitions; the counted text does not use them, and the span is taken whole so that the file version is unambiguous.\par\medskip
\noindent\textit{Editorial note.} The source URL is the source-level Stacks Project URL, as required when more than one card is produced from one source; the per-claim tag is carried by claim\_key and by this note. The counted statement and proof cite no bibliography entry, so this card needs no bibliography.\par\medskip
\noindent\textit{Allowed standard prerequisites (editorial).}\begin{itemize}
\item Relative pseudo-coherence of a complex is preserved by localizing the algebra at an element, which gives one of the two implications. invoked at proof lines 21683-21685. Disposes of the implication from the second to the first condition.
\item The algebra can be presented by a polynomial ring in the elements and their complementary coefficients modulo the relation expressing that the elements generate the unit ideal. used at proof lines 21686-21695. Produces the polynomial presentation to which the transfers of relative pseudo-coherence are applied.
\item The quotient presenting the algebra is pseudo-coherent as a module over the polynomial ring and remains so after inverting one of the elements, and relative pseudo-coherence transfers to a larger algebra under this hypothesis. used at proof lines 21695-21711. Moves the hypothesis from the localized algebras to the polynomial ring and back, giving the other implication.
\end{itemize}\medskip
\ifdefined\hypertarget\hypertarget{sourcepointer-1}{}\fi
\noindent\textit{Source reference.} Original lemma 0679: sources/stacks-more-algebra/pinned-original-source.tex, source lines 21499--21499. The reference is cited at proof lines 21683-21685 for the implication from relative pseudo-coherence to its localizations; the converse, which is the substance of the proof, is inside the counted span. The lemma itself is the subject of a separate card of this lane with disjoint coordinates.\par\medskip
\ifdefined\hypertarget\hypertarget{sourcepointer-2}{}\fi
\noindent\textit{Source reference.} Original lemma 064Z: sources/stacks-more-algebra/pinned-original-source.tex, source lines 16283--16283. The reference is applied twice, at proof lines 21699-21703 and at proof lines 21709-21711, to move relative pseudo-coherence between the algebra, its presentation and the polynomial ring; the presentation that makes it applicable is inside the counted span.\par\medskip
\ifdefined\hypertarget\hypertarget{sourcepointer-3}{}\fi
\noindent\textit{Source reference.} Original lemma 066D: sources/stacks-more-algebra/pinned-original-source.tex, source lines 16404--16404. The reference glues pseudo-coherence over the presentation from the localizations at the chosen elements, applied at proof lines 21704-21706; the localizations themselves are produced inside the counted span.\par\medskip
\begin{lemma}
\label{lemma-glue-relative-pseudo-coherent}
Let $R \to A$ be a finite type ring map.
Let $K^\bullet$ be a complex of $A$-modules.
Let $m \in \mathbf{Z}$.
Let $f_1, \ldots, f_r \in A$ generate the unit ideal.
The following are equivalent
\begin{enumerate}
\item each $K^\bullet \otimes_A A_{f_i}$ is
$m$-pseudo-coherent relative to $R$, and
\item $K^\bullet$ is $m$-pseudo-coherent relative to $R$.
\end{enumerate}
The same equivalence holds for pseudo-coherence relative to $R$.
\end{lemma}

\begin{proof}
The implication (2) $\Rightarrow$ (1) is in
Lemma \ref{lemma-localize-relative-pseudo-coherent}.
Assume (1). Write $1 = \sum f_ig_i$ in $A$.
Choose a surjection
$R[x_1, \ldots, x_n, y_1, \ldots, y_r, z_1, \ldots, z_r] \to A$.
such that $y_i$ maps to $f_i$ and $z_i$ maps to $g_i$. Then we
see that there exists a surjection
$$
P = R[x_1, \ldots, x_n, y_1, \ldots, y_r, z_1, \ldots, z_r]/(\sum y_iz_i - 1)
\longrightarrow
A.
$$
Note that $P$ is pseudo-coherent as an
$R[x_1, \ldots, x_n, y_1, \ldots, y_r, z_1, \ldots, z_r]$-module
and that $P[1/y_i]$ is pseudo-coherent as an
$R[x_1, \ldots, x_n, y_1, \ldots, y_r, z_1, \ldots, z_r, 1/y_i]$-module.
Hence by
Lemma \ref{lemma-finite-push-pseudo-coherent}
we see that
$K^\bullet \otimes_A A_{f_i}$ is an $m$-pseudo-coherent complex
of $P[1/y_i]$-modules for each $i$.
Thus by
Lemma \ref{lemma-glue-pseudo-coherent}
we see that $K^\bullet$ is pseudo-coherent as a complex of
$P$-modules, and
Lemma \ref{lemma-finite-push-pseudo-coherent}
shows that  $K^\bullet$ is pseudo-coherent as a complex of
$R[x_1, \ldots, x_n, y_1, \ldots, y_r, z_1, \ldots, z_r]$-modules.
\end{proof}

\end{document}
